\documentclass[11pt]{article}
\usepackage[margin=18mm,headheight=22pt]{geometry}
\usepackage{amsmath}
\usepackage{fontspec}
\usepackage{unicode-math}
\setmainfont{NewCM10}[Extension=.otf,UprightFont=*-Book,BoldFont=*-Bold,ItalicFont=*-BookItalic,BoldItalicFont=*-BoldItalic]
\setsansfont{NewCMSans10}[Extension=.otf,UprightFont=*-Book,BoldFont=*-Bold,ItalicFont=*-BookOblique,BoldItalicFont=*-BoldOblique]
\setmathfont{NewCMMath-Book.otf}
\usepackage{xcolor}
\definecolor{FAUblue}{HTML}{04316A}
\definecolor{Paperblue}{HTML}{EEF2F7}
\definecolor{EEIgold}{HTML}{CC9C42}
\usepackage{etoolbox,enumitem,booktabs,tabularx,array,needspace,microtype}
\usepackage[most]{tcolorbox}
\usepackage{titlesec,fancyhdr}
\usepackage[hidelinks]{hyperref}
\usepackage{draftwatermark}
\SetWatermarkText{Made by Thanate Amornjiraporn (Focus)}
\SetWatermarkColor{black!7}\SetWatermarkFontSize{30pt}\SetWatermarkAngle{45}
\setlength{\parindent}{0pt}
\setlength{\parskip}{5pt}
\setlist{itemsep=2pt,topsep=3pt,leftmargin=1.6em,beginpenalty=10000,endpenalty=10000}
\raggedbottom
\clubpenalty=10000 \widowpenalty=10000 \displaywidowpenalty=10000 \interlinepenalty=10000
\setcounter{secnumdepth}{0}
\setcounter{tocdepth}{1}
\pagestyle{fancy}\fancyhf{}
\fancyhead[L]{\small\sffamily\color{FAUblue} From Fundamentals to AI Apocalypse · APDR · Exercise 3}
\fancyhead[R]{\small\sffamily Computer Exercise 3: Worked solutions}
\fancyfoot[C]{\small\thepage}
\titleformat{\section}{\Large\bfseries\sffamily\color{FAUblue}}{}{0pt}{}
\titleformat{\subsection}{\large\bfseries\sffamily\color{FAUblue}}{}{0pt}{}
\titlespacing*{\section}{0pt}{16pt}{6pt}
\preto\section{\Needspace{8\baselineskip}}
\preto\subsection{\Needspace{6\baselineskip}}
\newtcolorbox{equationcard}[1]{colback=Paperblue,colframe=Paperblue,boxrule=0pt,arc=0pt,left=10pt,right=10pt,top=6pt,bottom=6pt,title={#1},fonttitle=\small\bfseries\sffamily,coltitle=FAUblue,before skip=8pt,after skip=8pt}
\newtcolorbox{workedbox}[1]{colback=white,colframe=EEIgold,boxrule=0pt,leftrule=3pt,arc=0pt,left=10pt,top=4pt,bottom=4pt,title={#1},fonttitle=\small\bfseries\sffamily,coltitle=FAUblue,attach title to upper={\par},before skip=8pt,after skip=8pt}
\newtcolorbox{recallbox}[1]{colback=white,colframe=FAUblue,boxrule=0pt,leftrule=3pt,arc=0pt,left=10pt,top=4pt,bottom=4pt,title={#1},fonttitle=\small\bfseries\sffamily,coltitle=FAUblue,attach title to upper={\par},before skip=8pt,after skip=8pt}
\usepackage{environ}
\newsavebox\measurebox
% Typeset the section once invisibly to measure it; start a new page if it will not fit in what is left.
\makeatletter
\newif\ifkeep@measuring
\NewEnviron{keepsection}{%
  \ifkeep@measuring \BODY \else
    \edef\keep@sec{\the\value{section}}\edef\keep@sub{\the\value{subsection}}%
    \setbox\measurebox\vbox{\hsize\linewidth\keep@measuringtrue\def\Needspace##1{}\BODY}%
    \setcounter{section}{\keep@sec}\setcounter{subsection}{\keep@sub}%
    \dimen2=\ht\measurebox \advance\dimen2\dp\measurebox
    \dimen@=\textheight
    \ifdim\pagegoal<\maxdimen
      \dimen@=\pagegoal \advance\dimen@-\pagetotal
      \ifdim\dimen2>\dimen@ \clearpage \dimen@=\textheight \fi
    \fi
    % if it fits, nothing inside may force an early break
    \ifdim\dimen2>\dimen@ \BODY \else \begingroup\def\Needspace##1{}\BODY\endgroup \fi
  \fi}
\makeatother
\newtcolorbox{notebox}{colback=Paperblue,colframe=Paperblue,boxrule=0pt,arc=0pt,left=8pt,right=8pt,top=4pt,bottom=4pt,before skip=6pt,after skip=6pt}
\newtcolorbox{plainbox}{colback=white,colframe=FAUblue,boxrule=0pt,leftrule=3pt,arc=0pt,left=10pt,right=8pt,top=4pt,bottom=4pt,before skip=6pt,after skip=8pt}
\usetikzlibrary{arrows.meta}
\title{\sffamily\bfseries\color{FAUblue}Computer Exercise 3: Worked solutions}
\author{Algorithms, Programming, and Data Representation \ensuremath{\cdot} Module 92313 \ensuremath{\cdot} Exercise 3\\[2pt]\small From Fundamentals to AI Apocalypse · personal study guide}
\date{}
\begin{document}
\maketitle
\begingroup\footnotesize\setlength{\parskip}{0pt}\tableofcontents\endgroup
\clearpage

\Needspace{9\baselineskip}
{\large Manual byte decoding, packet layouts, checksums, stream recovery, address translation, and round trips. Try each task before opening its worked steps. Personalized inputs and results are omitted; apply the methods to your own assigned data.}


\begin{keepsection}
\Needspace{5\baselineskip}
\section{Task 1: Decode an unsigned integer}

\textbf{Connected topic:} Related APDR guide section

\begin{keepsection}
\Needspace{6\baselineskip}
\subsection{Task 01 \ensuremath{\cdot} Decode an unsigned integer}

\Needspace{3\baselineskip}
Try first: Decode bytes manually in big-endian order, without int.from\_bytes.

\nopagebreak[4]{\small\itshape Answers: see “Task 1: Decode an unsigned integer” under Answers at the end.}
\end{keepsection}
\end{keepsection}

\begin{keepsection}
\Needspace{5\baselineskip}
\section{Task 2: Decode a signed two’s-complement integer}

\textbf{Connected topic:} Related APDR guide section

\begin{keepsection}
\Needspace{6\baselineskip}
\subsection{Task 02 \ensuremath{\cdot} Decode a signed two’s-complement integer}

\Needspace{3\baselineskip}
Try first: Decode the unsigned bit pattern, then interpret its top bit as the sign.

\nopagebreak[4]{\small\itshape Answers: see “Task 2: Decode a signed two’s-complement integer” under Answers at the end.}
\end{keepsection}
\end{keepsection}

\begin{keepsection}
\Needspace{5\baselineskip}
\section{Task 3: Decode IEEE 754 single precision}

\textbf{Connected topic:} Related APDR guide section

\begin{keepsection}
\Needspace{7\baselineskip}
\subsection{Task 03 \ensuremath{\cdot} Decode IEEE 754 single precision}

\Needspace{4\baselineskip}
Try first: Separate sign, exponent, and fraction bits, then apply the normal or subnormal formula.

\nopagebreak[4]{\small\itshape Answers: see “Task 3: Decode IEEE 754 single precision” under Answers at the end.}
\end{keepsection}
\end{keepsection}

\begin{keepsection}
\Needspace{5\baselineskip}
\section{Task 4: Infer byte order from a known range}

\textbf{Connected topic:} Related APDR guide section

\begin{keepsection}
\Needspace{7\baselineskip}
\subsection{Task 04 \ensuremath{\cdot} Infer byte order from a known range}

\Needspace{4\baselineskip}
Try first: Decode the same bytes both ways and use a plausible range to choose the unique match.

\nopagebreak[4]{\small\itshape Answers: see “Task 4: Infer byte order from a known range” under Answers at the end.}
\end{keepsection}
\end{keepsection}

\begin{keepsection}
\Needspace{5\baselineskip}
\section{Task 5: Parse a mixed packet layout}

\textbf{Connected topic:} Related APDR guide section

\begin{keepsection}
\Needspace{7\baselineskip}
\subsection{Task 05 \ensuremath{\cdot} Parse a mixed packet layout}

\Needspace{4\baselineskip}
Try first: Walk declared fields in order, consuming padding and decoding each field by its specified type.

\nopagebreak[4]{\small\itshape Answers: see “Task 5: Parse a mixed packet layout” under Answers at the end.}
\end{keepsection}
\end{keepsection}

\begin{keepsection}
\Needspace{5\baselineskip}
\section{Task 6: Validate an additive checksum}

\textbf{Connected topic:} Related APDR guide section

\begin{keepsection}
\Needspace{6\baselineskip}
\subsection{Task 06 \ensuremath{\cdot} Validate an additive checksum}

\Needspace{3\baselineskip}
Try first: Recompute a modulo-256 sum over payload bytes and compare the final byte.

\nopagebreak[4]{\small\itshape Answers: see “Task 6: Validate an additive checksum” under Answers at the end.}
\end{keepsection}
\end{keepsection}

\begin{keepsection}
\Needspace{5\baselineskip}
\section{Task 7: Parse a packet layout}

\textbf{Connected topic:} Related APDR guide section

\begin{keepsection}
\Needspace{7\baselineskip}
\subsection{Task 07 \ensuremath{\cdot} Parse a packet layout}

\Needspace{4\baselineskip}
Try first: Walk the declared fields in sequence, track the current byte offset, and decode by type.

\nopagebreak[4]{\small\itshape Answers: see “Task 7: Parse a packet layout” under Answers at the end.}
\end{keepsection}
\end{keepsection}

\begin{keepsection}
\Needspace{5\baselineskip}
\section{Task 8: Decode a stream and keep going after errors}

\textbf{Connected topic:} Related APDR guide section

\begin{keepsection}
\Needspace{7\baselineskip}
\subsection{Task 08 \ensuremath{\cdot} Decode a stream and keep going after errors}

\Needspace{4\baselineskip}
Try first: Process each packet independently, recording malformed packets rather than aborting the stream.

\nopagebreak[4]{\small\itshape Answers: see “Task 8: Decode a stream and keep going after errors” under Answers at the end.}
\end{keepsection}
\end{keepsection}

\begin{keepsection}
\Needspace{5\baselineskip}
\section{Task 9: Simulate first-fit heap allocations}

\textbf{Connected topic:} Related APDR guide section

\begin{keepsection}
\Needspace{7\baselineskip}
\subsection{Task 09 \ensuremath{\cdot} Simulate first-fit heap allocations}

\Needspace{4\baselineskip}
Try first: Allocate each request from the first free block large enough, returning its start and size or None.

\nopagebreak[4]{\small\itshape Answers: see “Task 9: Simulate first-fit heap allocations” under Answers at the end.}
\end{keepsection}
\end{keepsection}

\begin{keepsection}
\Needspace{5\baselineskip}
\section{Task 10: Check integer encode/decode round trips}

\textbf{Connected topic:} Related APDR guide section

\begin{keepsection}
\Needspace{6\baselineskip}
\subsection{Task 10 \ensuremath{\cdot} Check integer encode/decode round trips}

\Needspace{3\baselineskip}
Try first: Encode a value at a chosen width and verify decoding in each byte order returns it.

\nopagebreak[4]{\small\itshape Answers: see “Task 10: Check integer encode/decode round trips” under Answers at the end.}
\end{keepsection}
\end{keepsection}


\clearpage
\section{Answers}
\begin{keepsection}\subsection{Task 1: Decode an unsigned integer}

\begin{workedbox}{Task 01 \ensuremath{\cdot} worked solution}
\Needspace{12\baselineskip}
\begin{enumerate}
\item Start with value 0. For every byte, shift the accumulated value left by 8 bits, then OR in the new byte.
\item For bytes 0x12, 0x34 the calculation is (0 \ensuremath{\times} 256 + 18) \ensuremath{\times} 256 + 52 = 4660. This is positional notation in base 256.
\item Empty input naturally returns 0. Reject any endianness label other than big or little.
\item \Needspace{3\baselineskip}
\textbf{Reference implementation:}
\begin{tcolorbox}[colback=Paperblue,colframe=Paperblue,boxrule=0pt,arc=0pt,left=7pt,right=7pt,top=5pt,bottom=5pt]\small\sloppy
\noindent\hspace*{0em}\texttt{def bytes\_to\_uint(data, endianness="big"):}\par
\noindent\hspace*{1em}\texttt{if endianness == "big":}\par
\noindent\hspace*{3em}\texttt{value = 0}\par
\noindent\hspace*{3em}\texttt{for byte in data:}\par
\noindent\hspace*{5em}\texttt{value = (value \ensuremath{<}\ensuremath{<} 8) | byte}\par
\noindent\hspace*{3em}\texttt{return value}\par
\noindent\hspace*{1em}\texttt{if endianness == "little":}\par
\noindent\hspace*{3em}\texttt{return sum(byte \ensuremath{<}\ensuremath{<} (8 * i) for i, byte in enumerate(data))}\par
\noindent\hspace*{1em}\texttt{raise ValueError("endianness must be big or little")}\par
\end{tcolorbox}
\end{enumerate}
\end{workedbox}

\end{keepsection}
\begin{keepsection}\subsection{Task 2: Decode a signed two’s-complement integer}

\begin{workedbox}{Task 02 \ensuremath{\cdot} worked solution}
\Needspace{10\baselineskip}
\begin{enumerate}
\item For w bits, first compute the unsigned value u. If the high bit is 0, the signed value is u; otherwise subtract 2\textasciicircum{}w.
\item For the one-byte pattern 11111101, u = 253 and w = 8, so the value is 253 \ensuremath{-} 256 = \ensuremath{-}3.
\item The range is \ensuremath{-}2\textasciicircum{}(w\ensuremath{-}1) through 2\textasciicircum{}(w\ensuremath{-}1)\ensuremath{-}1. Check the all-zero and all-one patterns.
\item \Needspace{3\baselineskip}
\textbf{Reference implementation:}
\begin{tcolorbox}[colback=Paperblue,colframe=Paperblue,boxrule=0pt,arc=0pt,left=7pt,right=7pt,top=5pt,bottom=5pt]\small\sloppy
\noindent\hspace*{0em}\texttt{def bytes\_to\_int\_twos\_complement(data, endianness="big"):}\par
\noindent\hspace*{1em}\texttt{value = bytes\_to\_uint(data, endianness)}\par
\noindent\hspace*{1em}\texttt{width = 8 * len(data)}\par
\noindent\hspace*{1em}\texttt{if width and value \& (1 \ensuremath{<}\ensuremath{<} (width - 1)):}\par
\noindent\hspace*{3em}\texttt{value -= 1 \ensuremath{<}\ensuremath{<} width}\par
\noindent\hspace*{1em}\texttt{return value}\par
\end{tcolorbox}
\end{enumerate}
\end{workedbox}

\end{keepsection}
\begin{keepsection}\subsection{Task 3: Decode IEEE 754 single precision}

\begin{workedbox}{Task 03 \ensuremath{\cdot} worked solution}
\Needspace{15\baselineskip}
\begin{enumerate}
\item A binary32 word has 1 sign bit, 8 exponent bits, and 23 fraction bits. For a normal number, value = (\ensuremath{-}1)\textasciicircum{}s \ensuremath{\times} (1 + fraction/2\textasciicircum{}23) \ensuremath{\times} 2\textasciicircum{}(exponent\ensuremath{-}127).
\item If exponent is zero, use the subnormal significand fraction/2\textasciicircum{}23 and exponent \ensuremath{-}126, with no implicit leading one.
\item Exponent 255 represents infinity when fraction is zero and NaN otherwise. Require exactly 4 bytes and verify endianness first.
\item \Needspace{3\baselineskip}
\textbf{Reference implementation:}
\begin{tcolorbox}[colback=Paperblue,colframe=Paperblue,boxrule=0pt,arc=0pt,left=7pt,right=7pt,top=5pt,bottom=5pt]\small\sloppy
\noindent\hspace*{0em}\texttt{def bytes\_to\_float32(data, endianness="big"):}\par
\noindent\hspace*{1em}\texttt{if len(data) != 4:}\par
\noindent\hspace*{3em}\texttt{raise ValueError("float32 requires exactly 4 bytes")}\par
\noindent\hspace*{1em}\texttt{bits = bytes\_to\_uint(data, endianness)}\par
\noindent\hspace*{1em}\texttt{sign = -1 if bits \ensuremath{>}\ensuremath{>} 31 else 1}\par
\noindent\hspace*{1em}\texttt{exponent, fraction = (bits \ensuremath{>}\ensuremath{>} 23) \& 255, bits \& 0x7fffff}\par
\noindent\hspace*{1em}\texttt{if exponent == 255:}\par
\noindent\hspace*{3em}\texttt{return sign * float("inf") if fraction == 0 else float("nan")}\par
\noindent\hspace*{1em}\texttt{if exponent == 0:}\par
\noindent\hspace*{3em}\texttt{return sign * (fraction / 2**23) * 2**-126}\par
\noindent\hspace*{1em}\texttt{return sign * (1 + fraction / 2**23) * 2**(exponent - 127)}\par
\end{tcolorbox}
\end{enumerate}
\end{workedbox}

\end{keepsection}
\begin{keepsection}\subsection{Task 4: Infer byte order from a known range}

\begin{workedbox}{Task 04 \ensuremath{\cdot} worked solution}
\Needspace{11\baselineskip}
\begin{enumerate}
\item Compute big- and little-endian candidates independently.
\item A candidate matches when lower \ensuremath{\leq} value \ensuremath{\leq} upper. Return an order only if exactly one candidate matches; neither or both is ambiguous and should raise ValueError.
\item Test boundary values inclusively and include an ambiguous example.
\item \Needspace{3\baselineskip}
\textbf{Reference implementation:}
\begin{tcolorbox}[colback=Paperblue,colframe=Paperblue,boxrule=0pt,arc=0pt,left=7pt,right=7pt,top=5pt,bottom=5pt]\small\sloppy
\noindent\hspace*{0em}\texttt{def detect\_endianness(raw, known\_field\_range):}\par
\noindent\hspace*{1em}\texttt{lower, upper = known\_field\_range}\par
\noindent\hspace*{1em}\texttt{big = lower \ensuremath{<}= bytes\_to\_uint(raw, "big") \ensuremath{<}= upper}\par
\noindent\hspace*{1em}\texttt{little = lower \ensuremath{<}= bytes\_to\_uint(raw, "little") \ensuremath{<}= upper}\par
\noindent\hspace*{1em}\texttt{if big == little:}\par
\noindent\hspace*{3em}\texttt{raise ValueError("range does not identify unique byte order")}\par
\noindent\hspace*{1em}\texttt{return "big" if big else "little"}\par
\end{tcolorbox}
\end{enumerate}
\end{workedbox}

\end{keepsection}
\begin{keepsection}\subsection{Task 5: Parse a mixed packet layout}

\begin{workedbox}{Task 05 \ensuremath{\cdot} worked solution}
\Needspace{17\baselineskip}
\begin{enumerate}
\item Track a byte offset. For each field, slice exactly its width and advance the offset, including for padding.
\item Decode uint, int, float32, or bytes using that field’s endianness; padding produces no output key. Reject negative widths, truncated fields, and unknown types.
\item Check a layout with one padding byte and mixed-endian fields.
\item \Needspace{3\baselineskip}
\textbf{Reference implementation:}
\begin{tcolorbox}[colback=Paperblue,colframe=Paperblue,boxrule=0pt,arc=0pt,left=7pt,right=7pt,top=5pt,bottom=5pt]\small\sloppy
\noindent\hspace*{0em}\texttt{def parse\_packet(raw, layout):}\par
\noindent\hspace*{1em}\texttt{result, offset = \{\}, 0}\par
\noindent\hspace*{1em}\texttt{for field in layout:}\par
\noindent\hspace*{3em}\texttt{width = field["width"]}\par
\noindent\hspace*{3em}\texttt{end = offset + width}\par
\noindent\hspace*{3em}\texttt{if width \ensuremath{<} 0 or end \ensuremath{>} len(raw): raise ValueError("field exceeds packet")}\par
\noindent\hspace*{3em}\texttt{data, offset = raw[offset:end], end}\par
\noindent\hspace*{3em}\texttt{kind = field["type"]}\par
\noindent\hspace*{3em}\texttt{order = field.get("endianness", "big")}\par
\noindent\hspace*{3em}\texttt{if kind == "padding": continue}\par
\noindent\hspace*{3em}\texttt{if kind == "uint": value = bytes\_to\_uint(data, order)}\par
\noindent\hspace*{3em}\texttt{elif kind == "int": value = bytes\_to\_int\_twos\_complement(data, order)}\par
\noindent\hspace*{3em}\texttt{elif kind == "float32": value = bytes\_to\_float32(data, order)}\par
\noindent\hspace*{3em}\texttt{elif kind == "bytes": value = data}\par
\noindent\hspace*{3em}\texttt{else: raise ValueError(f"unsupported field type: \{kind\}")}\par
\noindent\hspace*{3em}\texttt{result[field["name"]] = value}\par
\noindent\hspace*{1em}\texttt{return result}\par
\end{tcolorbox}
\end{enumerate}
\end{workedbox}

\end{keepsection}
\begin{keepsection}\subsection{Task 6: Validate an additive checksum}

\begin{workedbox}{Task 06 \ensuremath{\cdot} worked solution}
\Needspace{10\baselineskip}
\begin{enumerate}
\item For packet bytes b\textsubscript{0} … b\textsubscript{n}, the expected checksum is sum(b\textsubscript{0} … b\textsubscript{n\ensuremath{-}1}) mod 256.
\item In Python, sum(raw[:-1]) \& 0xff keeps the low eight bits. Require a nonempty packet before indexing its last byte.
\item A checksum detects many accidental changes but is not authentication; distinct corruptions can collide.
\item \Needspace{3\baselineskip}
\textbf{Reference implementation:}
\begin{tcolorbox}[colback=Paperblue,colframe=Paperblue,boxrule=0pt,arc=0pt,left=7pt,right=7pt,top=5pt,bottom=5pt]\small\sloppy
\noindent\hspace*{0em}\texttt{def validate\_checksum(raw):}\par
\noindent\hspace*{1em}\texttt{return bool(raw) and (sum(raw[:-1]) \& 0xff) == raw[-1]}\par
\end{tcolorbox}
\end{enumerate}
\end{workedbox}

\end{keepsection}
\begin{keepsection}\subsection{Task 7: Parse a packet layout}

\begin{workedbox}{Task 07 \ensuremath{\cdot} worked solution}
\Needspace{8\baselineskip}
\begin{enumerate}
\item For each field, take raw[offset:offset + width] and advance the offset even for padding. Check that width is nonnegative and the slice stays within the packet.
\item Dispatch uint, int, float32, and bytes to their corresponding decoders. Padding consumes bytes without creating a result entry.
\item Reject unsupported types and malformed widths; test mixed endianness and padding with a small synthetic layout.
\end{enumerate}
\end{workedbox}

\end{keepsection}
\begin{keepsection}\subsection{Task 8: Decode a stream and keep going after errors}

\begin{workedbox}{Task 08 \ensuremath{\cdot} worked solution}
\Needspace{14\baselineskip}
\begin{enumerate}
\item Enumerate packets so error records retain the original packet index.
\item Check checksum before parsing. Wrap validation and parsing for one packet in try/except; append successful values to decoded and failures to errors.
\item A valid later packet must still be handled after a corrupt packet. Keep the original index in each error record.
\item \Needspace{3\baselineskip}
\textbf{Reference implementation:}
\begin{tcolorbox}[colback=Paperblue,colframe=Paperblue,boxrule=0pt,arc=0pt,left=7pt,right=7pt,top=5pt,bottom=5pt]\small\sloppy
\noindent\hspace*{0em}\texttt{class PageFaultError(Exception):}\par
\noindent\hspace*{1em}\texttt{pass}\par
\noindent\hspace*{0em}\texttt{}\par
\noindent\hspace*{0em}\texttt{class PageTableSim:}\par
\noindent\hspace*{1em}\texttt{def \_\_init\_\_(self, page\_size, page\_table):}\par
\noindent\hspace*{3em}\texttt{if page\_size \ensuremath{<}= 0: raise ValueError("page\_size must be positive")}\par
\noindent\hspace*{3em}\texttt{self.page\_size, self.page\_table = page\_size, page\_table}\par
\noindent\hspace*{1em}\texttt{def translate(self, virtual\_addr):}\par
\noindent\hspace*{3em}\texttt{page, offset = divmod(virtual\_addr, self.page\_size)}\par
\noindent\hspace*{3em}\texttt{if page not in self.page\_table: raise PageFaultError(page)}\par
\noindent\hspace*{3em}\texttt{return self.page\_table[page] * self.page\_size + offset}\par
\end{tcolorbox}
\end{enumerate}
\end{workedbox}

\end{keepsection}
\begin{keepsection}\subsection{Task 9: Simulate first-fit heap allocations}

\begin{workedbox}{Task 09 \ensuremath{\cdot} worked solution}
\Needspace{14\baselineskip}
\begin{enumerate}
\item Represent each free block as a start and length. Visit blocks in address order and choose the first with enough space.
\item Split the selected block into any unused prefix and suffix, or mark it fully allocated. Invalid negative requests should fail clearly.
\item Check that allocations do not overlap and that every successful range stays within the pool.
\item \Needspace{3\baselineskip}
\textbf{Reference implementation:}
\begin{tcolorbox}[colback=Paperblue,colframe=Paperblue,boxrule=0pt,arc=0pt,left=7pt,right=7pt,top=5pt,bottom=5pt]\small\sloppy
\noindent\hspace*{0em}\texttt{def simulate\_heap\_alloc(requests, pool\_size):}\par
\noindent\hspace*{1em}\texttt{if pool\_size \ensuremath{<} 0: raise ValueError("pool\_size cannot be negative")}\par
\noindent\hspace*{1em}\texttt{allocations, next\_address = [], 0}\par
\noindent\hspace*{1em}\texttt{for size in requests:}\par
\noindent\hspace*{3em}\texttt{if size \ensuremath{<} 0: raise ValueError("allocation size cannot be negative")}\par
\noindent\hspace*{3em}\texttt{if next\_address + size \ensuremath{<}= pool\_size:}\par
\noindent\hspace*{5em}\texttt{allocations.append((next\_address, size))}\par
\noindent\hspace*{5em}\texttt{next\_address += size}\par
\noindent\hspace*{3em}\texttt{else:}\par
\noindent\hspace*{5em}\texttt{allocations.append(None)}\par
\noindent\hspace*{1em}\texttt{return allocations}\par
\end{tcolorbox}
\end{enumerate}
\end{workedbox}

\end{keepsection}
\begin{keepsection}\subsection{Task 10: Check integer encode/decode round trips}

\begin{workedbox}{Task 10 \ensuremath{\cdot} worked solution}
\Needspace{14\baselineskip}
\begin{enumerate}
\item For an unsigned width w bytes, valid values satisfy 0 \ensuremath{\leq} value \ensuremath{<} 2\textasciicircum{}(8w). Validate width, value, and endianness before encoding.
\item Emit bytes from most significant to least significant for big-endian, and reverse that order for little-endian.
\item Round-trip zero, the maximum representable value, and a value with distinct bytes. The decoder should reproduce the original integer.
\item \Needspace{3\baselineskip}
\textbf{Reference implementation:}
\begin{tcolorbox}[colback=Paperblue,colframe=Paperblue,boxrule=0pt,arc=0pt,left=7pt,right=7pt,top=5pt,bottom=5pt]\small\sloppy
\noindent\hspace*{0em}\texttt{def uint\_to\_bytes(value, width, endianness="big"):}\par
\noindent\hspace*{1em}\texttt{if endianness not in ("big", "little") or width \ensuremath{<} 0:}\par
\noindent\hspace*{3em}\texttt{raise ValueError("invalid encoding parameters")}\par
\noindent\hspace*{1em}\texttt{if value \ensuremath{<} 0 or value \ensuremath{>}= 1 \ensuremath{<}\ensuremath{<} (8 * width):}\par
\noindent\hspace*{3em}\texttt{raise ValueError("value does not fit")}\par
\noindent\hspace*{1em}\texttt{shifts = range(width - 1, -1, -1) if endianness == "big" else range(width)}\par
\noindent\hspace*{1em}\texttt{return bytes((value \ensuremath{>}\ensuremath{>} (8 * shift)) \& 255 for shift in shifts)}\par
\end{tcolorbox}
\end{enumerate}
\end{workedbox}

\end{keepsection}

\vfill{\small Personal worked solutions based on the APDR Computer Exercise 3 sheet, explained in my own words. Examples are generic; use your own assigned inputs for any personalized result. Not an official FAU publication; check course details against current materials. Made by \textbf{Thanate Amornjiraporn (Focus)}.}

\end{document}
